% Gabarit pandoc des corrigés publiés (brut et guidé). Contenu sans marque ni nom de site.
% Compilation : lualatex (latexmk -lualatex).
% Le Markdown reste la source unique : ne pas éditer ce fichier .tex à la main.
\documentclass[11pt,a4paper]{article}

% --- Moteur, polices (Latin Modern) et langue ---------------------------------
\usepackage{iftex}
\ifLuaTeX\else\errmessage{Ce document se compile avec lualatex}\fi
\usepackage{amsmath,amssymb,mathtools}
\usepackage{fontspec}
% mathrm=sym : \mathrm prend ses glyphes dans la police mathématique (ℓ, μ, Ω), comme KaTeX.
\usepackage[math-style=TeX,bold-style=TeX,mathrm=sym]{unicode-math}
% Repli pour les glyphes absents de Latin Modern (indices et exposants Unicode,
% flèches, filets des schémas en texte, symboles) : DejaVu, puis Latin Modern Math
% (ℓ, lettres mathématiques Unicode des légendes).
\directlua{
  luaotfload.add_fallback("repliTexte", {"DejaVuSerif:mode=harf;", "DejaVuSans:mode=harf;", "file:latinmodern-math.otf:mode=harf;"})
  luaotfload.add_fallback("repliMono", {"DejaVuSansMono:mode=harf;", "DejaVuSans:mode=harf;", "file:latinmodern-math.otf:mode=harf;"})
}
\setmainfont{Latin Modern Roman}[RawFeature={fallback=repliTexte}]
\setsansfont{Latin Modern Sans}[RawFeature={fallback=repliTexte}]
\setmonofont{Latin Modern Mono}[RawFeature={fallback=repliMono}]
\setmathfont{Latin Modern Math}
\usepackage[french]{babel}
\frenchsetup{SmallCapsFigTabCaptions=false}

% --- Mise en page --------------------------------------------------------------
\usepackage[a4paper,margin=2.2cm,headheight=14pt,footskip=1.2cm]{geometry}
\usepackage{microtype}
\setlength{\emergencystretch}{4em}
\tolerance=1500
\setlength{\parindent}{0pt}
\setlength{\parskip}{0.55\baselineskip plus 2pt minus 1pt}
\usepackage{xcolor}
% Gris nommés : babel-french rend « ! » actif dans le document.
\definecolor{gris2}{gray}{0.98}
\definecolor{gris3}{gray}{0.97}
\definecolor{gris5}{gray}{0.95}
\definecolor{gris35}{gray}{0.65}
\definecolor{gris40}{gray}{0.60}
\definecolor{gris45}{gray}{0.55}
\definecolor{gris50}{gray}{0.50}
\definecolor{gris55}{gray}{0.45}
\definecolor{gris60}{gray}{0.40}
\definecolor{gris65}{gray}{0.35}
\definecolor{gris70}{gray}{0.30}
\definecolor{gris80}{gray}{0.20}
\definecolor{gris85}{gray}{0.15}
\usepackage{graphicx}
\usepackage{float}
\usepackage{caption}
\captionsetup{font=small,labelformat=empty,justification=centering}
\usepackage{booktabs,longtable,array,multirow,calc}
\usepackage{siunitx}
\sisetup{locale=FR}
\usepackage[version=4]{mhchem}
\usepackage{listings}
\usepackage{fancyhdr}
\usepackage{lastpage}
\usepackage[most]{tcolorbox}
\usepackage{enumitem}
\setlist{nosep,leftmargin=*,topsep=2pt}
\usepackage[hidelinks,unicode]{hyperref}
\usepackage{bookmark}
\hypersetup{pdftitle={ATS — Anglais 2026},pdfauthor={Rédigé par IA},pdfcreator={LaTeX},pdfsubject={Corrigé brut}}
\usepackage{xurl}

% --- Pandoc : listes, tableaux, figures ----------------------------------------
\providecommand{\passthrough}[1]{#1}
\providecommand{\tightlist}{\setlength{\itemsep}{0pt}\setlength{\parskip}{0pt}}

\makeatletter
\def\maxwidth{\ifdim\Gin@nat@width>\linewidth\linewidth\else\Gin@nat@width\fi}
\def\maxheight{\ifdim\Gin@nat@height>0.55\textheight 0.55\textheight\else\Gin@nat@height\fi}
\makeatother
\setkeys{Gin}{width=\maxwidth,height=\maxheight,keepaspectratio}
\setcounter{secnumdepth}{-\maxdimen}
\providecommand{\pandocbounded}[1]{#1}
\usepackage{longtable,booktabs,array}

% --- Code (listings) -----------------------------------------------------------
\lstset{
  basicstyle=\ttfamily\small,
  columns=fullflexible,
  keepspaces=true,
  breaklines=true,
  breakatwhitespace=false,
  postbreak=\mbox{\textcolor{gris50}{$\hookrightarrow$}\space},
  frame=single,
  rulecolor=\color{gris35},
  framesep=4pt,
  xleftmargin=5pt,xrightmargin=5pt,
  backgroundcolor=\color{gris3},
  keywordstyle=\bfseries,
  commentstyle=\color{gris65},
  stringstyle=\color{gris80},
  showstringspaces=false,
  tabsize=4,
  extendedchars=true,
  aboveskip=6pt,belowskip=6pt,
}

% --- Macros des corrigés (KaTeX ↔ LaTeX) ---------------------------------------
% Les corrigés sont écrits pour KaTeX : on définit ici ce que KaTeX connaît et
% que LaTeX n'a pas, sans jamais retoucher le Markdown.
\providecommand{\R}{\mathbb{R}}
\providecommand{\N}{\mathbb{N}}
\providecommand{\Z}{\mathbb{Z}}
\providecommand{\Q}{\mathbb{Q}}
\providecommand{\C}{\mathbb{C}}
\providecommand{\dd}{\mathrm{d}}
\providecommand{\llbracket}{\lBrack}
\providecommand{\rrbracket}{\rBrack}
\providecommand{\lBrace}{\lbrace}
\providecommand{\rBrace}{\rbrace}
\providecommand{\Box}{\square}
\providecommand{\Diamond}{\lozenge}
\providecommand{\varnothing}{\emptyset}
\providecommand{\argmax}{\operatorname*{arg\,max}}
\providecommand{\argmin}{\operatorname*{arg\,min}}
\renewcommand{\leq}{\leqslant}
\renewcommand{\geq}{\geqslant}

% Équation hors texte : réduite à la largeur de la ligne si elle déborde.
\newsavebox{\boiteequation}
\newcommand{\equationajustee}[1]{%
  \sbox{\boiteequation}{$\displaystyle #1$}%
  \ifdim\wd\boiteequation>\linewidth
    \[\resizebox{\linewidth}{!}{\usebox{\boiteequation}}\]%
  \else
    \[#1\]%
  \fi}

% --- Encadrés pédagogiques (sobres, lisibles en noir et blanc) -----------------
\tcbset{
  encadre/.style={
    enhanced,breakable,sharp corners,
    colback=white,colframe=gris55,
    boxrule=0.4pt,leftrule=2.2pt,
    left=6pt,right=6pt,top=3pt,bottom=3pt,
    before skip=6pt,after skip=6pt,
    fonttitle=\bfseries\small,coltitle=black,
    attach title to upper,after title={\quad},
    parbox=false,
  }
}
\newtcolorbox{indication}[1]{encadre,colframe=gris45,colback=gris2,title={Indication~#1}}
\newtcolorbox{methode}{encadre,colframe=gris70,title={Méthode}}
\newtcolorbox{piege}{encadre,colframe=gris85,boxrule=0.8pt,leftrule=2.2pt,title={Piège}}
\newtcolorbox{jury}{encadre,colframe=gris60,colback=gris5,title={Jury}}
\newtcolorbox{encadrenote}{encadre,colframe=gris40}

% --- Titres de section ---------------------------------------------------------
\usepackage{titlesec}
\titleformat{\section}{\Large\bfseries}{}{0pt}{}
\titleformat{\subsection}{\large\bfseries}{}{0pt}{}
\titleformat{\subsubsection}{\normalsize\bfseries}{}{0pt}{}
\titlespacing*{\section}{0pt}{14pt plus 4pt minus 2pt}{6pt}
\titlespacing*{\subsection}{0pt}{10pt plus 3pt minus 2pt}{4pt}
\titlespacing*{\subsubsection}{0pt}{8pt plus 2pt minus 2pt}{2pt}

% --- En-tête et pied de page ---------------------------------------------------
\pagestyle{fancy}
\fancyhf{}
\fancyhead[L]{\small\textsc{Corrigé brut} · ATS ATS}
\fancyhead[R]{\small Anglais 2026}
\fancyfoot[L]{\footnotesize Corrigé rédigé par IA, non officiel}
\fancyfoot[R]{\footnotesize \thepage\,/\,\pageref*{LastPage}}
\renewcommand{\headrulewidth}{0.4pt}
\renewcommand{\footrulewidth}{0pt}
\fancypagestyle{plain}{\fancyhf{}\fancyfoot[R]{\footnotesize \thepage\,/\,\pageref*{LastPage}}\fancyfoot[L]{\footnotesize Corrigé rédigé par IA, non officiel}\renewcommand{\headrulewidth}{0pt}}

\begin{document}
\thispagestyle{plain}
\begin{center}
  {\small\scshape Corrigé brut}\par\smallskip
  {\LARGE\bfseries ATS — Anglais 2026\par}\medskip
  {\large ATS · ATS · Anglais · 2026\par}
\end{center}
\vspace{2pt}\hrule\vspace{8pt}

\emph{Corrigé rédigé par IA, non officiel. Signalez toute erreur.}

\textbf{Problématique} : l'épreuve ne comporte aucune rédaction. Elle mesure la maîtrise automatisée de la grammaire et du lexique (100 items), puis la compréhension fine de quatre articles de presse (37 items). Chaque réponse est justifiée par la règle ou par la phrase du texte qui la fonde.

Durée : 2 h, sans document ni appareil. Barème de chaque item : réponse juste +3, réponse fausse −1, pas de réponse 0. Une seule réponse par item.

\hypertarget{partie-i-qcm-de-grammaire-et-de-lexique-questions-1-uxe0-100}{%
\section{Partie I --- QCM de grammaire et de lexique (questions 1 à 100)}\label{partie-i-qcm-de-grammaire-et-de-lexique-questions-1-uxe0-100}}

\begin{longtable}[]{@{}
  >{\raggedright\arraybackslash}p{(\columnwidth - 4\tabcolsep) * \real{0.0644}}
  >{\raggedright\arraybackslash}p{(\columnwidth - 4\tabcolsep) * \real{0.2292}}
  >{\raggedright\arraybackslash}p{(\columnwidth - 4\tabcolsep) * \real{0.6764}}@{}}
\toprule\noalign{}
\begin{minipage}[b]{\linewidth}\raggedright
N°
\end{minipage} & \begin{minipage}[b]{\linewidth}\raggedright
Réponse
\end{minipage} & \begin{minipage}[b]{\linewidth}\raggedright
Justification
\end{minipage} \\
\midrule\noalign{}
\endhead
\bottomrule\noalign{}
\endlastfoot
1 & c) \emph{talking} & \emph{enjoy} se construit avec la forme en -ing : \emph{We really enjoyed talking to him}. \\
2 & b) \emph{visiting} & \emph{be worth} + V-ing : \emph{Rome is well worth visiting} (« vaut la peine d'être visitée »). \\
3 & b) \emph{buy} & \emph{had better} + base verbale sans \emph{to} : \emph{We had better buy}. \\
4 & b) \emph{Given} & \emph{Given} (« étant donné, compte tenu de ») introduit un groupe nominal : \emph{Given the current global warming}. \\
5 & a) \emph{both} & Pour deux éléments pris ensemble, avec un nom et un verbe au pluriel : \emph{both days suit me}. \emph{either} et \emph{each} exigeraient le singulier (\emph{either day suits me}) ; \emph{every} ne s'emploie pas pour deux. \\
6 & a) \emph{tougher} & \emph{get tough on crime} = se montrer plus sévère envers la délinquance ; avec \emph{get} et un comparatif : \emph{get tougher on crime}. \\
7 & b) \emph{of} & \emph{disapprove of something} = désapprouver quelque chose. \\
8 & a) \emph{meeting} & \emph{to} est ici une préposition, suivie d'un nom ou d'une forme en -ing : \emph{I look forward to meeting her}. \\
9 & d) \emph{exploration} & Il faut un nom : \emph{Oil and gas exploration \ldots{} is getting more intense} (l'exploration pétrolière et gazière). \\
10 & b) \emph{working} & \emph{seem to be} + V-ing : \emph{Our website seems to be working normally again}. \\
11 & b) \emph{resume} & \emph{resume} (ou \emph{résumé}) est le mot américain pour le CV. \emph{CVs} est au pluriel alors que l'on s'adresse à un seul candidat ; \emph{pay slip} = fiche de paie ; \emph{bulletin} est un faux ami (bulletin d'information). \\
12 & c) \emph{forward} & \emph{put forward a suggestion} = avancer, proposer une suggestion. \\
13 & b) \emph{each other} & Pronom réciproque : \emph{Jack and Tom have known each other for many years}. \emph{themselves} serait réfléchi (chacun se connaît lui-même). \\
14 & c) \emph{to be} & \emph{consider something to be the driving force} (proposition infinitive après \emph{consider}). \\
15 & d) \emph{happens} & Après \emph{if} (condition portant sur l'avenir), présent simple et non \emph{will} : \emph{If the problem happens again, \ldots{} we'll be able}. \emph{comes out} (sortir, être publié) n'a pas le bon sens ; \emph{came} casserait la concordance avec \emph{we'll}. \\
16 & c) \emph{though} & \emph{It looks as though} (= \emph{as if}) \emph{it's going to be a splendid afternoon}. \\
17 & a) \emph{rise} & \emph{a pay rise} (britannique ; \emph{a raise} en américain). \emph{augmentation} (accroissement, au sens technique) ne s'emploie pas pour un salaire, et \emph{raises} est au pluriel après \emph{a}. \\
18 & a) \emph{hard} & \emph{a hard worker} = un gros travailleur. \emph{hardly} est un adverbe (« à peine ») ; \emph{indolent} et \emph{idle} (paresseux) contredisent \emph{conscientious}. \\
19 & b) \emph{fully} & \emph{a fully qualified engineer} = un ingénieur pleinement diplômé. \\
20 & b) \emph{no} & \emph{There are no actual pieces of evidence}. \emph{not} exigerait \emph{not any} ; \emph{none} est un pronom ; \emph{any} donnerait le sens inverse. \\
21 & b) \emph{into} & \emph{divided into two parts}. \\
22 & a) \emph{either} & Dans une phrase déjà négative (\emph{doesn't want}), « non plus » se dit \emph{either} en fin de phrase. \emph{neither} exigerait une inversion sans négation (\emph{neither does he}) ; \emph{as well} et \emph{also} s'emploient en phrase affirmative. \\
23 & c) \emph{promotional} & Adjectif épithète : \emph{a strong promotional strategy}. \\
24 & b) \emph{proposing} & Relative réduite au participe en -ing (sens actif) : \emph{Any government proposing tax cuts is bound to be popular}. \\
25 & d) \emph{outside} & \emph{just outside the city} = juste à l'extérieur de la ville. \emph{at the city} et \emph{on the city} ne se disent pas ; \emph{into} marque un mouvement. \\
26 & c) \emph{enough} & \emph{not enough chairs} (la salle était trop pleine : \emph{so popular that\ldots{}}). \\
27 & b) \emph{a little} & \emph{a little} + indénombrable = un peu (\emph{a few} pour les dénombrables). \emph{little} seul signifie « peu » (sens restrictif) ; dans une demande polie on emploie \emph{some}, pas \emph{any}. \\
28 & b) \emph{enforces} & \emph{enforce the law} = faire respecter, faire appliquer la loi. \emph{inflict} = infliger ; \emph{compells} est mal orthographié (\emph{compels}) et n'a pas ce sens. \\
29 & d) \emph{got up} & \emph{It's high time} + prétérit modal (valeur d'irréel) : \emph{It's high time we got up}. \\
30 & d) \emph{don't} & \emph{No, I don't (like fish)} : auxiliaire \emph{do} au présent, première personne. \\
31 & c) \emph{am} & Après \emph{when} temporel, présent et non futur : \emph{When I am older, I want to\ldots{}}. \\
32 & b) \emph{a} & \emph{become a football player} : article indéfini devant une profession ; \emph{a} devant une consonne. \\
33 & c) \emph{out of} & \emph{run out of food} = se trouver à court de nourriture (\emph{It was tough}). \\
34 & b) \emph{knowing} & Préposition + V-ing : \emph{Without knowing much about the situation}. \\
35 & a) \emph{who} & \emph{The person who invited me over} : \emph{who} sujet, antécédent humain. \emph{whom} est complément ; \emph{which} s'emploie pour les choses ; \emph{what} n'a pas d'antécédent. \\
36 & b) \emph{clockwork} & \emph{run like clockwork} = se dérouler comme sur des roulettes. \emph{a breeze} s'emploie dans \emph{be a breeze}, pas dans \emph{run like}. \\
37 & b) \emph{embarrassing} & La situation provoque la gêne : \emph{embarrassing} (-ing = qui cause). \emph{embarrassed} qualifie la personne gênée. \\
38 & c) \emph{was greeted} & Passif au prétérit : \emph{I was greeted by all my friends}. Le \emph{present perfect} est exclu avec un repère passé révolu. \\
39 & b) \emph{to carry out} & \emph{carry out an internship} = effectuer un stage (\emph{carry out} = réaliser, mener à bien). \\
40 & d) \emph{of} & \emph{tears of joy}. \\
41 & a) \emph{be allowed} & \emph{too young to be allowed in a pub} : infinitif passif ; après \emph{to}, base verbale (\emph{be}), d'où l'exclusion de \emph{being}. \\
42 & b) \emph{are being} & \emph{You are being naughty!} = tu n'es pas sage en ce moment (comportement temporaire). \\
43 & b) \emph{have you been doing} & \emph{lately} appelle le \emph{present perfect}, ici progressif (activité continue) : \emph{What have you been doing with your life lately?} \\
44 & a) \emph{badly} & \emph{do badly} = mal marcher (\emph{business is doing very badly}). \emph{hardly} = à peine ; \emph{harshly} = durement ; \emph{toughly} n'existe pas dans ce sens. \\
45 & a) \emph{out} & \emph{sort out a problem} = régler un problème. \\
46 & d) \emph{What} & \emph{What is going on here?} \\
47 & b) \emph{better} & \emph{You had better go now} = tu ferais mieux de partir. \\
48 & c) \emph{earlier / more} & \emph{The earlier you get up, the more you can do} : \emph{the} + comparatif, \emph{the} + comparatif. \emph{more earlier} est un double comparatif fautif. \\
49 & c) \emph{on} & \emph{what is on my mind} = ce qui me préoccupe, ce que j'ai en tête. \\
50 & c) \emph{himself} & \emph{think for oneself} : pronom réfléchi \emph{himself} (en un mot). \\
51 & c) \emph{needn't} & \emph{You needn't worry} = tu n'as pas besoin de t'inquiéter. \emph{don't have} exigerait \emph{to} ; \emph{must} et \emph{have to} donnent le sens inverse. \\
52 & d) \emph{latest} & \emph{his latest hit} = son tout dernier tube, le plus récent. \emph{last} signifie le dernier d'une série terminée. \\
53 & c) \emph{could} & \emph{if} + prétérit modal : \emph{if you could take a break}. \emph{can} et \emph{will} sont exclus ; \emph{might} n'a pas de sens de capacité ici. \\
54 & a) \emph{cake} & \emph{a piece of cake} = un jeu d'enfant. \\
55 & b) \emph{heard} & Prétérit simple pour l'événement ponctuel : \emph{we were reading \ldots{} when we heard a big noise}. \\
56 & c) \emph{fine} & \emph{a fine} = une amende. \emph{penalty} désigne une sanction en général ou un penalty ; \emph{reward} = récompense. \\
57 & a) \emph{about} & \emph{be passionate about something}. \\
58 & d) ∅ (aucun article) & Pluriel générique sans article : \emph{Cats are usually very independent animals}. \emph{Every} et \emph{Any} exigent le singulier ; \emph{The cats} désignerait des chats précis. \\
59 & b) \emph{most} & \emph{the most important soft skill} (superlatif : \emph{the most} + adjectif long). \\
60 & b) \emph{Electronics} & \emph{Electronics} (l'électronique), nom de discipline en -ics comme \emph{physics}, \emph{mathematics}. \emph{Electronic} est un adjectif. \\
61 & c) \emph{going} & Présent progressif à valeur de futur : \emph{Are you going out tonight?} \\
62 & b) \emph{Spaniard} & \emph{a Spaniard} = un Espagnol. \emph{Spanish} est un adjectif (ou désigne la langue) ; \emph{a Spanish} est fautif. \\
63 & a) \emph{walking distance} & \emph{within walking distance from/of my flat} = à distance de marche. \\
64 & c) \emph{within} & \emph{within two weeks} = en moins de deux semaines, dans un délai de deux semaines. \emph{less than} exigerait \emph{in less than}. \\
65 & c) \emph{through} & \emph{Monday through Friday} (américain ; \emph{Monday to Friday} en britannique). \\
66 & b) \emph{breaking up} & \emph{break up with someone} = rompre avec quelqu'un. \emph{break out} = éclater ; \emph{break into} = entrer par effraction. \\
67 & c) \emph{Whose} & Interrogatif possessif : \emph{Whose phone has she borrowed?} \emph{Who's} = \emph{who is / who has}. \\
68 & c) \emph{Although} & \emph{Although} + proposition. \emph{Despite} et \emph{In spite of} se construisent avec un nom ou une forme en -ing ; \emph{And yet} ne peut pas ouvrir une concessive placée avant la principale. \\
69 & b) so as to not to (pour \emph{so as not to}) & Le sens impose \emph{so as not to be disturbed} (« afin de ne pas être dérangés »). La seule option qui porte la négation est b), imprimée \emph{so as to not to} : le second \emph{to} est en trop, ce qui ressemble à une coquille d'énoncé. a) donne le contresens « afin d'être dérangés » ; \emph{for not to} est fautif ; \emph{with a view to} se construit avec V-ing. \\
70 & a) \emph{besides} & \emph{besides} = en plus de, outre. \emph{beside} = à côté de. \\
71 & b) \emph{across} & \emph{run across the road} = traverser la route en courant (d'un bord à l'autre d'une surface). \\
72 & c) \emph{above} & \emph{above the clouds} : au-dessus des nuages, le sol est caché. \\
73 & a) \emph{opposite} & \emph{opposite Joe} = en face de Joe. \\
74 & a) \emph{spend} & \emph{would rather} + base verbale : \emph{He would rather spend a holiday in Egypt than in Tunisia}. \\
75 & c) \emph{beside} & \emph{beside the church} = à côté de l'église. \emph{in front} exigerait \emph{of} (\emph{in front of}). \\
76 & b) \emph{Despite} & \emph{Despite the pandemic} (+ nom). \emph{In spite} exigerait \emph{of} ; \emph{Although} exige une proposition ; \emph{Because of} donne un contresens. \\
77 & b) \emph{have visited} & \emph{present perfect} avec une période qui inclut le présent : \emph{We have visited three Asian countries in the last two months}. \\
78 & a) \emph{of} & \emph{die of natural causes}. \emph{due} exigerait \emph{due to}. \\
79 & c) \emph{has she} & \emph{Neither has she} : inversion, et reprise de l'auxiliaire \emph{have} du \emph{present perfect}, à la 3e personne. \\
80 & a) \emph{far} & \emph{far better} = bien meilleur (comme \emph{much better}). \\
81 & b) \emph{series} & \emph{a series of accidents} : \emph{series} est invariable (singulier et pluriel identiques). \\
82 & a) \emph{much} & \emph{information} est indénombrable : \emph{as much information as we can}. \\
83 & c) \emph{should be} & \emph{should} exprime une probabilité forte : \emph{he should be here in a few minutes}. \emph{ought to being} est fautif (\emph{ought to be}). \\
84 & b) \emph{success} & \emph{I wish him success} (\emph{wish somebody something}). \\
85 & a) \emph{would get} & \emph{if you were more organized} (irréel du présent) → \emph{would} + base verbale : \emph{You would get better results}. \\
86 & c) \emph{better known as} & Participe passé en position de relative réduite : \emph{The part of Northern California better known as Silicon Valley is very famous}. a) ajouterait un second verbe conjugué. L'option c) reprend \emph{as}, déjà imprimé dans la phrase : coquille sans incidence sur le choix. \\
87 & d) \emph{would have reacted} & \emph{She wondered how she would have reacted ten years before} : irréel du passé (\emph{would have} + participe passé). \\
88 & b) \emph{likely} & \emph{it's very likely (that) it will snow} : adjectif. \emph{probably}, \emph{possibly} et \emph{certainly} sont des adverbes. \\
89 & d) \emph{could} & \emph{hundreds of lives could be saved} (modal + infinitif passif). \emph{could have} exigerait \emph{been saved} ; \emph{were able} et \emph{were allowed to} sont mal construits. \\
90 & a) \emph{must} & \emph{you must make an effort} (obligation). \emph{needn't} et \emph{can't} donnent un contresens. \\
91 & c) \emph{checking} & \emph{After checking the results} (préposition + V-ing). \emph{having check} est fautif (\emph{having checked}). \\
92 & c) \emph{a little} & \emph{a little more than his wife}. \emph{few}, \emph{many} et \emph{a few} portent sur des noms dénombrables. \\
93 & b) \emph{must have} & \emph{somebody must have told him} : \emph{must have} + participe passé = déduction sur le passé. \emph{can't have} donne l'inverse. \\
94 & b) \emph{inviting} & \emph{What about inviting your parents?} (\emph{What about} + V-ing). \\
95 & b) \emph{was she} & Inversion sujet-auxiliaire après \emph{Not only} en tête : \emph{Not only was she happy\ldots, but she was also grateful}. \\
96 & a) \emph{to survive} & \emph{it was hard to survive} (\emph{it is} + adjectif + infinitif complet). \\
97 & d) \emph{have} & \emph{why I have to pay} : le blanc est suivi de \emph{to}, ce qui exclut les modaux \emph{must}, \emph{needn't}, \emph{should} (jamais suivis de \emph{to}). \\
98 & a) ∅ (pas de relatif) & Relatif complément d'objet, donc omissible : \emph{the man she met at my wedding}. \emph{whose} et \emph{what} sont exclus, \emph{than} est un comparatif. \\
99 & c) \emph{few} & \emph{very few people} (dénombrable). \emph{little} et \emph{much} portent sur des indénombrables ; \emph{many} contredit le froid. \\
100 & d) \emph{any smarter} & \emph{he is not any smarter than you are} (\emph{not any} + comparatif). \emph{as smart} exigerait \emph{as you are} ; \emph{as smarter} est fautif. \\
\end{longtable}

\hypertarget{partie-ii-text-one-solo-travel-questions-101-uxe0-110}{%
\section{Partie II --- Text one : solo travel (questions 101 à 110)}\label{partie-ii-text-one-solo-travel-questions-101-uxe0-110}}

\begin{longtable}[]{@{}
  >{\raggedright\arraybackslash}p{(\columnwidth - 4\tabcolsep) * \real{0.0642}}
  >{\raggedright\arraybackslash}p{(\columnwidth - 4\tabcolsep) * \real{0.3359}}
  >{\raggedright\arraybackslash}p{(\columnwidth - 4\tabcolsep) * \real{0.5699}}@{}}
\toprule\noalign{}
\begin{minipage}[b]{\linewidth}\raggedright
N°
\end{minipage} & \begin{minipage}[b]{\linewidth}\raggedright
Réponse
\end{minipage} & \begin{minipage}[b]{\linewidth}\raggedright
Justification
\end{minipage} \\
\midrule\noalign{}
\endhead
\bottomrule\noalign{}
\endlastfoot
101 & c) \emph{a play} & \emph{The Unfriend} est une comédie du West End et le texte la reprend par \emph{The play may be funny}. \\
102 & c) \emph{people see it as a way to make long-lasting friendships} & \emph{the chance to make lifelong friendships is at the heart of it} : \emph{lifelong} = \emph{long-lasting}. d) (\emph{make friends on the way}) est trop faible : le texte insiste sur des amitiés durables. \\
103 & c) \emph{she got lots of support from the other group members} & \emph{The support you can get from a group of women you don't even know is phenomenal.} La composition du groupe (16 femmes) était connue d'avance : ce n'est pas une découverte. \\
104 & c) \emph{she had never had so much fun before} & \emph{{[}I had{]} never laughed so hard in my whole life} = elle ne s'était jamais autant amusée. a) (\emph{been laughed at}) et b) (\emph{made fun of people}) sont des contresens. \\
105 & a) \emph{she put up an Australian friend} & \emph{host a friend} = héberger un ami (\emph{put somebody up}). Rien n'indique une location ou des hôtes payants. \\
106 & b) \emph{had the feeling that sharing a room could have been enjoyable} & \emph{missing out} = passer à côté de quelque chose ; \emph{slightly} atténue le regret, d'où \emph{This time I'll share a room}. c) (\emph{really badly regretted}) et d) (\emph{lonely and depressed}) exagèrent. \\
107 & a) \emph{Jordan} & Elle est allée en Jordanie ; \emph{Norway is next} (voyage à venir) ; l'amie australienne est venue chez elle. \\
108 & c) \emph{may be shy and need help from professional tour guides to socialize with others} & G Adventures conçoit des circuits qui laissent le temps de se rencontrer et forme ses guides à créer de la camaraderie : les clients ont besoin d'aide pour se lier. Les autres options prêtent à Marceau des injonctions absentes du texte. \\
109 & c) \emph{identify who was not on good terms with other customers} & \emph{to spot when people might not be getting on} = repérer ceux qui ne s'entendent pas. Le texte ne parle ni de bagarre ni de rapport au supérieur. \\
110 & b) \emph{an unexpected and pleasant experience with his room mate} & Il redoutait \emph{a dreadful experience} mais s'est lié au point de rester ami avec son compagnon de chambre de 75 ans : expérience inattendue et agréable. \\
\end{longtable}

\hypertarget{partie-iii-text-two-les-pandas-de-hong-kong-questions-111-uxe0-121}{%
\section{Partie III --- Text two : les pandas de Hong Kong (questions 111 à 121)}\label{partie-iii-text-two-les-pandas-de-hong-kong-questions-111-uxe0-121}}

\begin{longtable}[]{@{}
  >{\raggedright\arraybackslash}p{(\columnwidth - 4\tabcolsep) * \real{0.1055}}
  >{\raggedright\arraybackslash}p{(\columnwidth - 4\tabcolsep) * \real{0.2197}}
  >{\raggedright\arraybackslash}p{(\columnwidth - 4\tabcolsep) * \real{0.6448}}@{}}
\toprule\noalign{}
\begin{minipage}[b]{\linewidth}\raggedright
N°
\end{minipage} & \begin{minipage}[b]{\linewidth}\raggedright
Réponse
\end{minipage} & \begin{minipage}[b]{\linewidth}\raggedright
Justification
\end{minipage} \\
\midrule\noalign{}
\endhead
\bottomrule\noalign{}
\endlastfoot
111 & a) \emph{beloved} & \emph{the beloved bears} : \emph{beloved} = adoré, chéri (adjectif). \emph{loving} = affectueux (qui aime). \\
112 & b) \emph{on} & \emph{capitalize on something} = tirer parti de quelque chose. \\
113 & a) \emph{struggled} & \emph{has struggled to attract tourists} = a eu du mal à attirer les touristes. \\
114 & a) \emph{crackdown} & \emph{a crackdown on dissent} = une répression de la contestation. \\
115 & b) \emph{currently} & \emph{Hong Kong currently has six pandas}. \\
116 & b) \emph{breeding} & \emph{breeding efforts} = tentatives de reproduction. \emph{breading} (paner) est un piège orthographique. \\
117 & c) \emph{cubs} & \emph{The cubs} : les petits (au pluriel, ce sont des jumeaux). \\
118 & a) \emph{welcomed} & \emph{In September, Hong Kong also welcomed a pair of pandas} (prétérit). \emph{has welcome} est fautif. \\
119 & c) \emph{officials} & \emph{Hong Kong officials} = les responsables, les autorités. \emph{politics} désigne la politique, pas les personnes. \\
120 & d) \emph{bears} & \emph{the two Beijing-gifted bears} : le panda géant est un ours. Ils ne sont ni des bébés ni des animaux de compagnie. \\
121 & a) \emph{struggling} & \emph{has been struggling financially} = connaît des difficultés financières (\emph{present perfect} progressif). \\
\end{longtable}

\hypertarget{partie-iv-text-three-how-to-stop-an-argument-questions-122-uxe0-126}{%
\section{Partie IV --- Text three : How to Stop an Argument? (questions 122 à 126)}\label{partie-iv-text-three-how-to-stop-an-argument-questions-122-uxe0-126}}

\begin{longtable}[]{@{}
  >{\raggedright\arraybackslash}p{(\columnwidth - 4\tabcolsep) * \real{0.0904}}
  >{\raggedright\arraybackslash}p{(\columnwidth - 4\tabcolsep) * \real{0.3500}}
  >{\raggedright\arraybackslash}p{(\columnwidth - 4\tabcolsep) * \real{0.5296}}@{}}
\toprule\noalign{}
\begin{minipage}[b]{\linewidth}\raggedright
N°
\end{minipage} & \begin{minipage}[b]{\linewidth}\raggedright
Réponse
\end{minipage} & \begin{minipage}[b]{\linewidth}\raggedright
Justification
\end{minipage} \\
\midrule\noalign{}
\endhead
\bottomrule\noalign{}
\endlastfoot
122 & b) \emph{full of} & \emph{are littered with scoreboards} = sont jonchés de tableaux de score, donc \emph{full of}. \\
123 & a) \emph{count} & \emph{meaningful relationships} = des relations qui comptent. \emph{humdrum} = banal. \\
124 & a) \emph{sticks to their argument and doesn't budge} & \emph{digging further into their entrenched position} = s'enfoncer dans sa position, ne pas bouger d'un pouce. \\
125 & a) \emph{to do better than the other} & \emph{to outdo the other} = surpasser l'autre. \\
126 & b) \emph{agree with each other} & \emph{we should be on the same page} = nous devrions être d'accord. \\
\end{longtable}

\hypertarget{partie-v-text-four-my-week-with-a-man-bag-questions-127-uxe0-137}{%
\section{Partie V --- Text four : My week with a man bag (questions 127 à 137)}\label{partie-v-text-four-my-week-with-a-man-bag-questions-127-uxe0-137}}

\begin{longtable}[]{@{}
  >{\raggedright\arraybackslash}p{(\columnwidth - 4\tabcolsep) * \real{0.0727}}
  >{\raggedright\arraybackslash}p{(\columnwidth - 4\tabcolsep) * \real{0.3250}}
  >{\raggedright\arraybackslash}p{(\columnwidth - 4\tabcolsep) * \real{0.5723}}@{}}
\toprule\noalign{}
\begin{minipage}[b]{\linewidth}\raggedright
N°
\end{minipage} & \begin{minipage}[b]{\linewidth}\raggedright
Réponse
\end{minipage} & \begin{minipage}[b]{\linewidth}\raggedright
Justification
\end{minipage} \\
\midrule\noalign{}
\endhead
\bottomrule\noalign{}
\endlastfoot
127 & b) \emph{the catching on of ``ladies' bags'' for men} & L'article porte sur la mode des sacs à main féminins chez les hommes (\emph{fashionable men have begun to appear in public wearing \ldots{} ladies' handbags}). \emph{catch on} = devenir à la mode. \\
128 & d) \emph{she is probably surprised at seeing a man using a handbag} & \emph{Perhaps admiring is the wrong word} et \emph{So, is this, like, a thing now?} traduisent la surprise, non l'admiration. \\
129 & d) \emph{is unaware of the situation} & Il découvre une tendance qu'il ignorait (\emph{Again, I have missed something. How can I hope to catch up?}). \\
130 & d) \emph{pricey handbags from glamorous brands} & Chalamet porte un petit sac Chanel, Pharrell Williams a relancé le Speedy de Louis Vuitton : des sacs coûteux de grandes marques. Le sac Gucci à anse en bambou est celui de Harry Styles. \\
131 & c) \emph{is totally lost} & \emph{I am absolutely at sea} = je suis complètement perdu ; la suite le confirme (\emph{They all look the same to me}). \\
132 & c) \emph{finds it difficult to see differences between bags} & \emph{They all look the same to me} : il ne distingue pas les sacs entre eux. \\
133 & b) \emph{is what the man is looking for} & \emph{a shearling bag \ldots{} that seems to fit the bill} = qui correspond à ce qu'il cherche. \\
134 & d) \emph{doesn't like her husband sporting a handbag} & \emph{pointing to the shearling bag over my shoulder} : c'est le sac porté par son mari qui la dérange. \\
135 & a) \emph{embarrassed and self-conscious} & \emph{The extent of the discomfort I feel}, \emph{undiluted mortification} : il est gêné et mal à l'aise. \\
136 & c) \emph{don't care about what people use for a bag} & \emph{no one cares what you are putting your shopping in at the self-checkout} : personne ne juge. \\
137 & b) \emph{outdated} & \emph{obsolete} = dépassé, périmé (\emph{outdated}). \\
\end{longtable}

\hypertarget{grille-ruxe9capitulative}{%
\section{Grille récapitulative}\label{grille-ruxe9capitulative}}

\begin{longtable}[]{@{}
  >{\raggedright\arraybackslash}p{(\columnwidth - 2\tabcolsep) * \real{0.5000}}
  >{\raggedright\arraybackslash}p{(\columnwidth - 2\tabcolsep) * \real{0.5000}}@{}}
\toprule\noalign{}
\begin{minipage}[b]{\linewidth}\raggedright
Items
\end{minipage} & \begin{minipage}[b]{\linewidth}\raggedright
Réponses
\end{minipage} \\
\midrule\noalign{}
\endhead
\bottomrule\noalign{}
\endlastfoot
1--10 & 1 c · 2 b · 3 b · 4 b · 5 a · 6 a · 7 b · 8 a · 9 d · 10 b \\
11--20 & 11 b · 12 c · 13 b · 14 c · 15 d · 16 c · 17 a · 18 a · 19 b · 20 b \\
21--30 & 21 b · 22 a · 23 c · 24 b · 25 d · 26 c · 27 b · 28 b · 29 d · 30 d \\
31--40 & 31 c · 32 b · 33 c · 34 b · 35 a · 36 b · 37 b · 38 c · 39 b · 40 d \\
41--50 & 41 a · 42 b · 43 b · 44 a · 45 a · 46 d · 47 b · 48 c · 49 c · 50 c \\
51--60 & 51 c · 52 d · 53 c · 54 a · 55 b · 56 c · 57 a · 58 d · 59 b · 60 b \\
61--70 & 61 c · 62 b · 63 a · 64 c · 65 c · 66 b · 67 c · 68 c · 69 b · 70 a \\
71--80 & 71 b · 72 c · 73 a · 74 a · 75 c · 76 b · 77 b · 78 a · 79 c · 80 a \\
81--90 & 81 b · 82 a · 83 c · 84 b · 85 a · 86 c · 87 d · 88 b · 89 d · 90 a \\
91--100 & 91 c · 92 c · 93 b · 94 b · 95 b · 96 a · 97 d · 98 a · 99 c · 100 d \\
101--110 & 101 c · 102 c · 103 c · 104 c · 105 a · 106 b · 107 a · 108 c · 109 c · 110 b \\
111--120 & 111 a · 112 b · 113 a · 114 a · 115 b · 116 b · 117 c · 118 a · 119 c · 120 d \\
121--130 & 121 a · 122 b · 123 a · 124 a · 125 a · 126 b · 127 b · 128 d · 129 d · 130 d \\
131--137 & 131 c · 132 c · 133 b · 134 d · 135 a · 136 c · 137 b \\
\end{longtable}

\end{document}
